\chapter[Band 23: Ungerichtete Graphen]{Band 23 auf einen Blick}
\noindent\textbf{Ungerichtete Graphen}\par\medskip
Die Kanten bleiben geordnete Paare. Ein eigenes Symmetrieaxiom
ergänzt den gerichteten Graphbegriff; Schleifenfreiheit und Zusammenhang
werden hier noch nicht gefordert.

\begin{BandUebersicht}
\textbf{Symmetrische Graphstruktur}
\UebFormel{\begin{aligned}
&\UndirectedGraphStruct{V,E}\ \leftrightarrow\\
&\quad\GraphStruct{V,E}\land\\
&\quad\forall x,y\in V\\
&\qquad((x,y)\in E\to(y,x)\in E).
\end{aligned}}
\UebQuellen{\FormulaRefAuto{UndirectedGraphDef}[def]}
&
\textbf{Adjazenz ist richtungsunabhängig}
\UebFormel{\begin{aligned}
&\UndirectedGraphStruct{V,E},\
 x,y\in V\ \vdash\\
&\qquad (x,y)\in E\ \leftrightarrow\ (y,x)\in E.
\end{aligned}}
\UebQuellen{\FormulaRefAuto{UndirectedAdjacencySymmetric}[thm]} \\
\addlinespace[.8ex]

\textbf{Ungerichtete Nachbarschaft}
\UebFormel{\begin{aligned}
&\UndirectedGraphStruct{V,E},\ x\in V\ \vdash\\
&\GraphNeighbors E x
  \coloneqq\{y\in V\mid(x,y)\in E\}.
\end{aligned}}
\UebQuellen{\FormulaRefAuto{UndirectedNeighborhoodDef}[def]}
&
\textbf{Gegenseitige Nachbarschaft}
\UebFormel{\begin{aligned}
&\UndirectedGraphStruct{V,E},\
 x,y\in V\ \vdash\\
&y\in\GraphNeighbors E x
  \ \leftrightarrow\ x\in\GraphNeighbors E y.
\end{aligned}}
\UebQuellen{\FormulaRefAuto{UndirectedNeighborhoodSymmetric}[thm]} \\
\addlinespace[.8ex]

\textbf{Induzierter Teilgraph}
\UebFormel{\begin{aligned}
&U\subseteq V,\qquad
 \DirectedInducedEdges E U=E\cap(U\times U).
\end{aligned}}
\UebQuellen{\FormulaRefAuto{DirectedInducedEdgeRelationDef}[def]}
&
\textbf{Symmetrie bleibt erhalten}
\UebFormel{\begin{aligned}
&\UndirectedGraphStruct{V,E},\ U\subseteq V\
 \vdash\\
&\qquad\UndirectedGraphStruct{U,\DirectedInducedEdges E U}.
\end{aligned}}
\UebQuellen{\FormulaRefAuto{UndirectedInducedSubgraph}[thm]} \\
\addlinespace[.8ex]
\end{BandUebersicht}
